\scanpage{11}\mspage{10}

Par suite, si $\delta$ est non dégénéré i.e. $\chi(\delta)\ne0$, alors $\varphi(\delta)$ est non dégénérée, \deleted{et si $\delta$ est une \unsure{polar|polarisation}} \inserted{et si $\delta$ est une} polarisation, $-\varphi(\delta)\inserted{=\delta'}$ est une polarisation. On \inserted{déduit} d'ailleurs de la relation plus haut, \deleted{\unsure{que|puisque}} $\deg\widetilde\delta=\chi(\delta)^2$, $\deg\widetilde{\delta'}=\chi(\delta')^2$, $\deg\widetilde\delta\,\widetilde{\delta'}=\chi(\delta)^2\chi(\delta')^2$,

\[
\delta'=-\varphi(\delta)\qquad\underset{\chi(\delta)^{2\unsure{n|g}}}{\shortparallel}
\]

% slanted
\margin{$\delta'=-\varphi(\delta)$ à priori, on a égalité des carrés ; mais les deux côtés sont des fonctions polynomiales, sur $\mathrm{NS}$ positives si $\delta$ est une polarisation}

\[
\chi(\deleted{\unsure{\varphi(}}\delta'\deleted{\unsure{)}})=\chi(\delta)^{n-1}
\]

si $\deleted{\unsure{\varphi(}}\delta)$ est non dégénéré, \deleted{alors $\delta$ \unsure{l'est}} $\chi(\delta)$ l'est.

Notons aussi que si $\delta'=-\varphi(\delta)$, $\delta''=-\varphi(\delta')$, on a

\[
\widetilde{\delta'}=\delta^{-1}\chi(\delta),\qquad
\deleted{\varphi(\varphi(\delta))=}\delta''{}^{-1}=\delta'{}^{-1}\chi(\delta')=\delta\chi(\delta)^{-1}\chi(\delta')=\delta\chi(\delta)^{n-2}
\]

\[
\text{donc}\qquad\varphi(\varphi(\delta))=\delta\chi(\delta)^{n-2}
\]

Si \deleted{\unsure{$g$}} $n=1$, i.e. $X$ une courbe elliptique, on trouve \deleted{$\varphi(\delta)={}$ polarisation canonique,} $\delta\chi(\delta)$ i.e. $\varphi(\varphi(\delta))$ est la polarisation canonique ; \deleted{[NB $\varphi($} mieux, pour tt $\delta$, $\varphi(\delta)$ est la polarisation canonique. Si $n\geq2$, la formule $\varphi(\varphi(\delta))=\delta\chi(\delta)^{n-2}$ inspire confiance directement\ldots.

\uline{NB} Il faudrait prouver que sur le \deleted{\unsure{schéma|schéma de}} \deleted{\unsure{topos}} \inserted{site} modulaire des schémas abéliens polarisés $(A,\delta)$ \inserted{en car. 0}, avec degré de polarisation donné $\chi(\delta)$, la polarisation $\varphi(\delta)=\delta'$ \deleted{\unsure{ind}} sur $B$ n'est plus divisible, \deleted{\unsure{ce que l'hom|i.e. que \illegible{1}}} ce qui montrerait que $\varphi(\delta)$ est ce qu'on peut faire de mieux dans ce genre comme ``inversion'' d'une polarisation. \deleted{\unsure{donc|On}} \inserted{En fait,} on s'attend à ce que $\varphi(\delta)$ engendre le groupe \deleted{par \unsure{regarder au pt de vue}} des sections de $\underline{\mathrm{NS}}_{B/S}$ sur une quelconque composante connexe du \deleted{\unsure{site modulaire}} \deleted{topos} \inserted{site}, modulaire\ldots
